\documentclass[11pt]{article}
\usepackage[margin=1in]{geometry}
\usepackage{amsmath,amssymb,amsthm,hyperref}
\newtheorem{theorem}{Theorem}
\title{A two-point Gibbs family freezes at the minimum}
\author{ziangni-sys}
\date{Conjecture 00000008486}
\begin{document}
\sloppy
\maketitle
\noindent Both source versions specify the Gibbs density
$\mu_t\propto e^{-f/t}$ as $t\downarrow0$, and assert that if the maximum
set of $f$ is a periodic orbit, then the limit is the uniform measure on
that orbit. The negative sign in the stated density makes this assertion
false even for a compact two-point system.

\begin{theorem}
There is a continuous function on a compact dynamical system whose maximum
set is a fixed orbit, but whose stated Gibbs family converges weakly to a
different point mass, rather than the uniform measure on the maximum orbit.
\end{theorem}
\begin{proof}
Take $\Omega=\{0,1\}$ with the discrete topology and its Borel sigma-algebra,
and let $T$ be the identity. Set $f(0)=0$ and $f(1)=1$. The space is compact,
both $T$ and $f$ are continuous, and the entire maximum set of $f$ is
$\{1\}$, a periodic orbit of period one. Its uniform measure is $\delta_1$.

Use counting measure on $\Omega$ as the reference measure for the Gibbs
density. For $t>0$, its partition function and normalized measure are
\[
 Z_t=\sum_{x\in\Omega}e^{-f(x)/t}=1+e^{-1/t},\qquad
 \mu_t=\frac{1}{1+e^{-1/t}}\delta_0+
       \frac{e^{-1/t}}{1+e^{-1/t}}\delta_1.
\]
Both weights are positive and sum to one, so this is an actual probability
measure with exactly the density prescribed in the source. Using the uniform
reference probability instead gives the same normalized Gibbs measure.

As $t\downarrow0$, the reciprocal $1/t$ tends to infinity, hence
$e^{-1/t}\to0$. Therefore the weights of $0$ and $1$ tend respectively to
$1$ and $0$. For every bounded continuous real function $g$ on $\Omega$,
\[
 \int g\,d\mu_t
 =\frac{g(0)}{1+e^{-1/t}}+
   \frac{e^{-1/t}g(1)}{1+e^{-1/t}}
 \longrightarrow g(0)=\int g\,d\delta_0.
\]
Thus $\mu_t$ converges weakly to $\delta_0$. Since
$\delta_0(\{0\})=1$ and $\delta_1(\{0\})=0$, these two probability
measures are distinct. The weak topology on probabilities on this finite
discrete space is Hausdorff, so $\mu_t$ cannot also converge to $\delta_1$.
This contradicts the asserted maximum-orbit conclusion.
\end{proof}

\section*{Scope}
The sign is exactly the negative sign appearing in both originals; it is
not replaced by the positive sign that would favor maxima. The reference
measure gives equal positive weight to both points, and the Gibbs measure
is explicitly normalized. No limit point, normalization, or orbit property
is assumed. Refuting the maximum-orbit clause suffices to disprove the
conjunction, without resolving the other equilibrium or oscillation claims.

\newpage
\raggedright
\section*{Actual probability measures and weak topology in Lean}
The state type is \texttt{Fin 2} with its discrete topology and actual Borel
measurable structure. \texttt{compact\_space}, \texttt{continuous\_T} and
\texttt{continuous\_f} prove the topological hypotheses. The definition
\texttt{maximizers} quantifies over every state; \texttt{maximum\_set}
proves it equals $\{1\}$. The theorem \texttt{maximum\_orbit} proves that
the range of all actual iterates of $T$ at $1$ is precisely this maximum set;
\texttt{maximum\_fixed} proves period one directly.

The nonnegative real weights \texttt{p t} and \texttt{q t} are defined using
the actual real exponential. Their sum is proved to be one.
The definition \texttt{raw t} is the sum of their scalar multiples of the
actual Dirac measures. Its total mass is proved to be one, giving the
bundled \texttt{ProbabilityMeasure} called \texttt{gibbs t}.
The theorem \texttt{partition\_function} evaluates the actual finite
exponential sum. For every state, \texttt{gibbs\_mass} proves exactly
\[
 \mu_t(\{x\})=\frac{e^{-f(x)/t}}{\sum_{y\in\Omega} e^{-f(y)/t}}.
\]
Point masses specify the whole measure on this finite space, so this is a
complete correspondence with the source's Gibbs family relative to counting
measure, rather than an asserted probability table.

The theorem \texttt{integral\_gibbs} computes genuine Bochner integrals
against the constructed measure. \texttt{weight\_limit}, \texttt{p\_limit}
and \texttt{q\_limit} prove the exponential and normalized weight limits.
The theorem \texttt{weak\_limit\_minimum} then proves convergence in
Mathlib's actual weak topology on \texttt{ProbabilityMeasure}, using its
characterization by integrals against every bounded continuous real function.
No custom convergence relation is substituted.

The definitions \texttt{point 0} and \texttt{point 1} are actual Dirac
probability measures. \texttt{point\_masses\_distinct} distinguishes them
by their measures of $\{0\}$. Finally, \texttt{not\_weak\_limit\_maximum}
negates actual weak convergence to \texttt{point 1}, using uniqueness of
limits in the Hausdorff probability-measure topology.

All convergence statements use the actual right-hand neighborhood filter
\texttt{nhdsWithin 0 (Set.Ioi 0)}. Totalized real division lets the definitions
also exist at temperature zero, but that value plays no role in these
positive-temperature limits.

\section*{Reproduction}
Lean 4.19.0 and Mathlib are pinned at commit
\begin{center}\small\texttt{c44e0c8ee63ca166450922a373c7409c5d26b00b}.\end{center}
From \texttt{lean/}, run \texttt{lake update}, \texttt{lake exe cache get},
and \texttt{lake build}; then check the complete source with
\begin{center}\small\texttt{lake env lean Main.lean -DwarningAsError=true}.\end{center}
The source prints five principal dependency audits. Only standard logical
axioms are used, with no admitted statements or auxiliary computations.
\end{document}
